% Generative Modelling
% Intended insertion point: after Architectures: Graph Neural Networks and before Inference and Model Compression.
% Uses standard packages already expected by this textbook: amsmath, amssymb, booktabs, tabularx, array, tikz, float.
% TikZ libraries used: arrows.meta, positioning, calc, fit, backgrounds.
% Assumes the textbook callout environments:
% calloutmemoryhook, calloutintuition, calloutinductivebias, calloutdegreeofexperiment,
% calloutmemorise, calloutexample, calloutpracticaltip, calloutspecialcase.

\chapter{Generative Modelling}
\label{chap:generative-modelling}

A generative model learns a mechanism that turns an easy source of randomness into samples distributed like the observed data:
\[
\underbrace{z\sim p_{\text{simple}}}_{\text{easy to sample}}
\quad\xrightarrow{\text{learned mechanism}}\quad
\underbrace{x\sim p_{\text{data}}}_{\text{complex target distribution}}.
\]
The useful distinction between generative families is therefore not simply their network architecture, but \textbf{what object is learned between the simple and data distributions}.
A U-Net or Transformer can parameterise several of these objects, so architecture, objective and sampler are separate design layers.
Autoregressive models learn conditional next-element distributions and were already developed for language generation in Section~13.7; this section focuses on latent-variable, adversarial, diffusion/score and flow-based generation.

\begin{table}[H]
\centering
\small
\begin{tabularx}{\linewidth}{>{\bfseries}l X X}
\toprule
Family & What is learned & Sampling path \\
\midrule
VAE & probabilistic latent encoder/decoder & sample latent $\rightarrow$ decode \\
GAN & direct generator trained against a discriminator/critic & noise $\rightarrow$ generator \\
Diffusion / score & denoising or score-related field across noise levels & noise $\rightarrow$ repeated reverse dynamics \\
Flow matching & time-dependent velocity field & noise $\rightarrow$ integrate an ODE \\
Autoregressive & conditional next-element distribution & repeatedly sample the next element \\
\bottomrule
\end{tabularx}
\caption{Generative families differ mainly in the learned object and the path used at inference; the neural backbone is a separate choice.}
\label{tab:generative-families}
\end{table}

\begin{calloutmemoryhook}
Generative models differ mainly in what they learn between simple probability and the data distribution.
\end{calloutmemoryhook}

\section{GANs: learn a generator through an opponent}

A \textbf{generative adversarial network (GAN)} trains a generator to produce samples that a discriminator cannot distinguish from observed data.
The two data paths are
\[
\underbrace{z}_{\text{random source}}
\rightarrow
\underbrace{G_{\theta_{\text{generator}}}(z)}_{\text{generated sample}}
\rightarrow
\underbrace{D_{\theta_{\text{discriminator}}}}_{\text{real / fake judgement}},
\qquad
\underbrace{x_{\text{real}}}_{\text{observed sample}}
\rightarrow
D_{\theta_{\text{discriminator}}}.
\]
The discriminator therefore supplies a \emph{learned loss surface}: it becomes better at exposing unrealistic samples while the generator becomes better at exploiting the current discriminator.
A classical objective is
\[
\min_{\theta_{\text{generator}}}\max_{\theta_{\text{discriminator}}}
\left[
\underbrace{\mathbb E_{x\sim p_{\text{data}}}\log D(x)}_{\text{reward recognising real data}}
+
\underbrace{\mathbb E_{z\sim p(z)}\log\!\left(1-D(G(z))\right)}_{\text{reward rejecting generated data}}
\right].
\]
This gives direct one-pass sampling after training, but optimisation is a coupled game rather than minimisation of one fixed supervised target.
If the generator finds samples that fool the current discriminator while ignoring some real modes, \textbf{mode collapse} can produce high-quality examples with poor distribution coverage.
Wasserstein-style objectives change the critic geometry to provide a more useful training signal, but do not remove the general cost of adversarially balancing two learned systems.

\begin{figure}[H]
\centering
\begin{tikzpicture}[
    >=Latex,
    node distance=9mm and 13mm,
    box/.style={draw, rounded corners=2pt, minimum width=23mm, minimum height=8mm, align=center, inner sep=3pt},
    smallbox/.style={draw, rounded corners=2pt, minimum width=18mm, minimum height=7mm, align=center, inner sep=2pt},
    arrow/.style={->, thick}
]
\node[box] (noise) {$z\sim p(z)$};
\node[box, right=of noise] (gen) {generator\\$G_\theta$};
\node[box, right=of gen] (fake) {$x_{\text{fake}}$};
\node[box, right=of fake] (disc) {discriminator\\$D_\phi$};
\node[smallbox, above=8mm of fake] (real) {$x_{\text{real}}$};
\node[smallbox, right=of disc] (judgement) {real / fake};
\draw[arrow] (noise) -- (gen);
\draw[arrow] (gen) -- (fake);
\draw[arrow] (fake) -- (disc);
\draw[arrow] (real) -| (disc);
\draw[arrow] (disc) -- (judgement);
\draw[->, thick, dashed, bend right=32] (disc.south) to node[below, align=center]{training signal\\changes with critic} (gen.south);
\end{tikzpicture}
\caption{GAN training is a feedback system: the discriminator learns to expose current generator errors, so the generator's training signal changes as its opponent changes.}
\label{fig:gan-flow}
\end{figure}

\begin{calloutinductivebias}
If a learned real/fake comparison provides a useful signal for perceptual sample quality, adversarial training can learn a strong direct generator; the same coupling can be restrictive when stable optimisation or complete mode coverage matters more.
\end{calloutinductivebias}

\begin{calloutmemoryhook}
Generator fools; discriminator exposes.
\end{calloutmemoryhook}

\section{VAEs: make the latent space sampleable}

Section~5.6 introduced the VAE objective; the architectural consequence is easier to see as a data-flow constraint.
A plain autoencoder only needs latent codes that reconstruct observed inputs, whereas a \textbf{variational autoencoder (VAE)} also forces those codes into a distribution from which new latent points can be sampled:
\[
\underbrace{x}_{\text{data}}
\rightarrow
\underbrace{q_{\theta_{\text{encoder}}}(z\mid x)}_{\text{distribution over latent codes}}
\rightarrow
\underbrace{z}_{\text{sampleable latent}}
\rightarrow
\underbrace{p_{\theta_{\text{decoder}}}(x\mid z)}_{\text{probabilistic decoder}}
\rightarrow
\hat x.
\]
The negative evidence lower bound used in Section~5.6 is already the whole trade-off:
\[
\underbrace{\mathcal L_{\text{VAE}}}_{\text{negative ELBO}}
=
\underbrace{-\mathbb E_{q(z\mid x)}\log p(x\mid z)}_{\text{reconstruct / explain observed data}}
+
\underbrace{D_{\mathrm{KL}}\!\left(q(z\mid x)\Vert p(z)\right)}_{\text{keep latent distribution compatible with the prior}}.
\]
The first term rewards detailed reconstruction; the second prevents the encoder from placing every input in an arbitrary isolated region that would be difficult to sample later.
The resulting tension is useful rather than accidental: stronger prior regularisation usually makes the latent space easier to sample but can remove information needed for reconstruction.

Sampling inside the encoder would ordinarily interrupt pathwise differentiation.
The \textbf{reparameterisation trick} moves the randomness into an external noise variable:
\[
\underbrace{z}_{\text{stochastic latent}}
=
\underbrace{\mu_{\theta_{\text{encoder}}}(x)}_{\text{latent centre}}
+
\underbrace{\sigma_{\theta_{\text{encoder}}}(x)}_{\text{latent scale}}
\odot
\underbrace{\epsilon}_{\epsilon\sim\mathcal N(0,I)}.
\]
Gradients can now flow through $\mu$ and $\sigma$ while $\epsilon$ supplies the stochasticity.

\begin{figure}[H]
\centering
\begin{tikzpicture}[
    >=Latex,
    node distance=9mm and 13mm,
    box/.style={draw, rounded corners=2pt, minimum width=25mm, minimum height=8mm, align=center, inner sep=3pt},
    arrow/.style={->, thick}
]
\node[box] (x) {data $x$};
\node[box, right=of x] (enc) {encoder\\$\mu(x),\sigma(x)$};
\node[box, right=of enc] (z) {$z=\mu+\sigma\odot\epsilon$};
\node[box, right=of z] (dec) {decoder\\$p(x\mid z)$};
\node[box, right=of dec] (xh) {reconstruction\\$\hat x$};
\node[box, below=9mm of z] (prior) {prior\\$p(z)$};
\draw[arrow] (x) -- (enc);
\draw[arrow] (enc) -- (z);
\draw[arrow] (z) -- (dec);
\draw[arrow] (dec) -- (xh);
\draw[<->, thick, dashed] (z) -- node[right, align=left]{KL keeps encoded\\latents sampleable} (prior);
\end{tikzpicture}
\caption{A VAE adds a distributional constraint to an autoencoder: latent codes must reconstruct the data while remaining compatible with a simple sampling prior.}
\label{fig:vae-flow}
\end{figure}

\begin{calloutmemoryhook}
A VAE reconstructs well without letting the latent space stop being sampleable.
\end{calloutmemoryhook}

\section{The score function: local geometry of a distribution}

For diffusion, the useful object is often not the density value itself but how that density changes around the current sample.
The \textbf{score function} at noise level $t$ is
\[
\boxed{
\underbrace{s_t(x)}_{\text{score field}}
=
\underbrace{\nabla_x\log p_t(x)}_{\text{direction of increasing log-density}}
}.
\]
The score is a vector field over sample space: near a high-density region it points in directions that locally increase probability density.
This does not require the normalising constant of $p_t(x)$ because differentiation removes constants independent of $x$.
Do not confuse this with the likelihood-ratio / policy-gradient ``score-function estimator'', which differentiates $\log p_\theta(x)$ with respect to parameters rather than the sample coordinates.

For Gaussian corruption
\[
\underbrace{x_t}_{\text{noisy sample}}
=
\underbrace{\alpha_t x_0}_{\text{remaining signal}}
+
\underbrace{\sigma_t\epsilon}_{\text{noise}},
\qquad \epsilon\sim\mathcal N(0,I),
\]
the conditional score is
\[
\underbrace{\nabla_{x_t}\log p(x_t\mid x_0)}_{\text{score of the corruption kernel}}
=
\underbrace{-\frac{x_t-\alpha_t x_0}{\sigma_t^2}}_{-\epsilon/\sigma_t}.
\]
This is the bridge between \textbf{score prediction} and the common diffusion objective that predicts the added noise $\epsilon$: they are scaled versions of the same local information under Gaussian corruption.

\begin{figure}[H]
\centering
\begin{tikzpicture}[scale=1.0, >=Latex]
\draw[->] (-3.2,0) -- (3.2,0) node[right] {$x_1$};
\draw[->] (0,-2.2) -- (0,2.2) node[above] {$x_2$};
\draw[thick] (-1.45,0) ellipse (0.85 and 1.15);
\draw[thick] (1.45,0) ellipse (0.85 and 1.15);
\fill (-1.45,0) circle (2pt);
\fill (1.45,0) circle (2pt);
\foreach \x/\y/\tx/\ty in {
-2.8/1.5/-1.8/0.7,
-2.9/-1.2/-1.8/-0.5,
-1.8/1.8/-1.55/0.8,
-1.7/-1.8/-1.5/-0.8,
-0.4/1.5/-1.1/0.6,
-0.4/-1.4/-1.05/-0.6,
2.8/1.5/1.8/0.7,
2.9/-1.2/1.8/-0.5,
1.8/1.8/1.55/0.8,
1.7/-1.8/1.5/-0.8,
0.4/1.5/1.1/0.6,
0.4/-1.4/1.05/-0.6}
{\draw[->, thick] (\x,\y) -- (\tx,\ty);}
\node[align=center] at (-1.45,-2.0) {high-density\\mode};
\node[align=center] at (1.45,-2.0) {high-density\\mode};
\end{tikzpicture}
\caption{Schematic score field for a bimodal distribution. Arrows indicate local directions of increasing log-density; the exact field between overlapping modes depends on their relative densities and scales.}
\label{fig:score-field}
\end{figure}

\begin{calloutmemoryhook}
Score points uphill in log-density.
\end{calloutmemoryhook}

\section{Diffusion: design corruption, learn the reverse}

\paragraph{Forward process.}
A diffusion model first defines a corruption process that gradually removes information from data until the terminal distribution is close to simple Gaussian noise.
For a standard discrete Gaussian process,
\[
\underbrace{q(x_t\mid x_{t-1})}_{\text{one corruption step}}
=
\mathcal N\!\left(
\underbrace{\sqrt{1-\beta_t}\,x_{t-1}}_{\text{retain signal}},
\underbrace{\beta_t I}_{\text{add noise}}
\right).
\]
Writing $\alpha_t=1-\beta_t$ and $\bar\alpha_t=\prod_{s=1}^{t}\alpha_s$ collapses all earlier corruption steps into one direct sample:
\[
\boxed{
\underbrace{x_t}_{\text{sample at noise level }t}
=
\underbrace{\sqrt{\bar\alpha_t}\,x_0}_{\text{remaining data signal}}
+
\underbrace{\sqrt{1-\bar\alpha_t}\,\epsilon}_{\text{known Gaussian corruption}},
\qquad \epsilon\sim\mathcal N(0,I)
}.
\]
Training therefore does not need to simulate every previous corruption step before constructing $x_t$.
The schedule $\{\beta_t\}$ controls how signal-to-noise ratio is distributed across time and is a real research choice rather than bookkeeping.

\begin{figure}[H]
\centering
\begin{tikzpicture}[
    >=Latex,
    stage/.style={draw, rounded corners=2pt, minimum width=29mm, minimum height=27mm, inner sep=2pt},
    arrow/.style={->, thick}
]
\node[stage] (s0) at (0,0) {};
\node[stage] (s1) at (3.6,0) {};
\node[stage] (s2) at (7.2,0) {};
\node[stage] (s3) at (10.8,0) {};
% structured clusters
\foreach \x/\y in {-0.8/0.6,-0.65/0.45,-0.7/0.8,-0.5/0.65,0.7/0.6,0.55/0.45,0.75/0.8,0.5/0.7,-0.8/-0.65,-0.6/-0.5,-0.7/-0.8,-0.5/-0.65,0.7/-0.65,0.55/-0.5,0.75/-0.8,0.5/-0.7}
{\fill ($(s0.center)+(\x,\y)$) circle (1.35pt);}
% light corruption
\foreach \x/\y in {-0.95/0.85,-0.55/0.25,-0.9/0.35,-0.35/0.95,0.95/0.8,0.35/0.25,0.95/0.35,0.25/0.95,-0.95/-0.8,-0.35/-0.25,-0.9/-0.35,-0.25/-1.0,0.9/-0.85,0.35/-0.2,1.0/-0.4,0.25/-0.95}
{\fill ($(s1.center)+(\x,\y)$) circle (1.35pt);}
% heavy corruption
\foreach \x/\y in {-1.0/0.9,-0.7/0.15,-0.8/-0.4,-0.3/1.0,-0.1/0.35,0.25/0.9,0.45/-0.1,0.9/0.75,1.05/0.0,0.75/-0.8,0.2/-0.95,-0.3/-0.6,-0.9/-0.95,0.55/0.4,0.1/0.05,-0.5/0.55}
{\fill ($(s2.center)+(\x,\y)$) circle (1.35pt);}
% approximate isotropic noise
\foreach \x/\y in {-1.05/0.45,-0.8/-0.65,-0.55/0.95,-0.45/0.15,-0.25/-0.25,-0.1/0.7,0.05/-0.95,0.2/0.2,0.35/0.85,0.45/-0.55,0.65/0.4,0.8/-0.15,0.95/0.75,1.0/-0.8,-0.75/0.1,0.1/-0.35}
{\fill ($(s3.center)+(\x,\y)$) circle (1.35pt);}
\draw[arrow] (s0) -- node[above] {add noise} (s1);
\draw[arrow] (s1) -- node[above] {add noise} (s2);
\draw[arrow] (s2) -- node[above] {add noise} (s3);
\node[below=2mm of s0] {$x_0$: structured data};
\node[below=2mm of s1] {$x_t$: partial signal};
\node[below=2mm of s2] {$x_t$: low SNR};
\node[below=2mm of s3] {$x_T\approx\mathcal N(0,I)$};
\draw[<-, very thick, dashed] ($(s0.north)+(0,3mm)$) -- node[above]{learned reverse dynamics} ($(s3.north)+(0,3mm)$);
\end{tikzpicture}
\caption{Diffusion separates the problem into a known forward corruption path and a learned reverse path. The point clouds are schematic: increasing $t$ destroys multimodal structure until only an approximately simple noise distribution remains.}
\label{fig:diffusion-forward-reverse}
\end{figure}

\paragraph{Reverse process.}
Generation starts at the easy end, $x_T\sim\mathcal N(0,I)$, and repeatedly predicts enough local information to move toward cleaner samples.
In a DDPM-style parameterisation,
\[
\underbrace{p_\theta(x_{t-1}\mid x_t)}_{\text{learned reverse step}}
=
\mathcal N\!\left(
\underbrace{\mu_\theta(x_t,t)}_{\text{predicted denoising move}},
\underbrace{\sigma_t^2 I}_{\text{reverse-step randomness}}
\right).
\]
A common network target is the exact noise $\epsilon$ used to construct $x_t$, because that target is available during training.
From the same noisy input $x_t$, related parameterisations can instead predict the clean sample $x_0$, the score, or a velocity-like $v$ target; these are different numerical views of closely related reverse information rather than separate generative families.

\begin{table}[H]
\centering
\small
\begin{tabularx}{\linewidth}{>{\bfseries}l X X}
\toprule
Prediction target & Network learns & Main consequence \\
\midrule
$\epsilon$-prediction & added Gaussian noise & simple supervised target; common DDPM parameterisation \\
$x_0$-prediction & clean sample estimate & directly interpretable reconstruction target; scaling changes with SNR \\
score prediction & $\nabla_x\log p_t(x)$ & exposes the reverse-SDE connection directly \\
$v$-prediction & a weighted combination of clean signal and noise & changes target scaling across noise levels; common in modern latent systems \\
\bottomrule
\end{tabularx}
\caption{Diffusion prediction parameterisation is a design choice inside the same broad reverse-process idea.}
\label{tab:diffusion-targets}
\end{table}

\begin{calloutmemoryhook}
Forward diffusion is designed corruption; the model learns the reverse information.
\end{calloutmemoryhook}

\section{DDPM and DDIM: separate the model from the sampler}

The denoising network and the numerical path used to turn noise into data are separate design layers.
A \textbf{DDPM} follows stochastic reverse transitions, injecting randomness while stepping through the learned reverse chain.
A canonical \textbf{DDIM} sampler can reuse a DDPM-trained network while following a non-Markovian path that becomes deterministic when its stochasticity parameter is set to zero.
This matters because changing the sampler can alter inference speed and trajectories without retraining the underlying denoiser.

\begin{table}[H]
\centering
\small
\begin{tabularx}{\linewidth}{>{\bfseries}l X X}
\toprule
 & DDPM & DDIM \\
\midrule
Learned denoiser & can be the same trained network & can be the same trained network \\
Reverse path & stochastic Markov chain & non-Markovian construction; can be deterministic \\
Noise during sampling & normally retained & optional; zero gives deterministic sampling \\
Practical use & canonical probabilistic reverse process & change trajectory / use fewer selected timesteps without changing the training family \\
\bottomrule
\end{tabularx}
\caption{DDPM versus DDIM is primarily a sampler distinction, not a statement that two different neural architectures must be trained.}
\label{tab:ddpm-ddim}
\end{table}

\begin{calloutmemoryhook}
Same denoiser; different sampling path.
\end{calloutmemoryhook}

\section{Continuous diffusion: SDE and probability-flow views}

The discrete sequence $x_0,x_1,\ldots,x_T$ can be taken to continuous time, where the forward corruption becomes a stochastic differential equation (SDE):
\[
\boxed{
\underbrace{dx}_{\text{state change}}
=
\underbrace{f(x,t)\,dt}_{\text{deterministic drift}}
+
\underbrace{g(t)\,dW_t}_{\text{random diffusion}}
}.
\]
The corresponding reverse-time SDE depends on the score of the noisy marginal distribution:
\[
\underbrace{dx}_{\text{reverse-time change}}
=
\underbrace{\left[f(x,t)-g(t)^2\nabla_x\log p_t(x)\right]dt}_{\text{drift corrected by the learned score}}
+
\underbrace{g(t)\,d\bar W_t}_{\text{reverse-time randomness}},
\qquad t:T\rightarrow0.
\]
The score therefore supplies exactly the missing information required to reverse a known diffusion process.
Under an exact score, the same sequence of marginal distributions can also be generated by the deterministic \textbf{probability-flow ODE}:
\[
\boxed{
\frac{dx}{dt}
=
\underbrace{f(x,t)-\frac{1}{2}g(t)^2\nabla_x\log p_t(x)}_{\text{deterministic velocity with the same marginals}}
}.
\]
This is the key bridge from diffusion to ODE-based generative modelling: stochastic and deterministic trajectories can represent the same evolving marginal densities while moving individual samples differently.

\begin{figure}[H]
\centering
\begin{tikzpicture}[
    >=Latex,
    node distance=10mm and 18mm,
    box/.style={draw, rounded corners=2pt, minimum width=28mm, minimum height=9mm, align=center, inner sep=3pt},
    arrow/.style={->, thick}
]
\node[box] (data) {data distribution\\$p_0$};
\node[box, right=45mm of data] (noise) {simple noise\\$p_T$};
\draw[arrow, bend left=18] (data) to node[above]{forward SDE: known corruption} (noise);
\draw[arrow, bend left=18] (noise) to node[below]{reverse SDE: score + randomness} (data);
\node[box, below=17mm of $(data)!0.5!(noise)$] (ode) {probability-flow ODE\\score + deterministic integration};
\draw[arrow] (noise) -- (ode);
\draw[arrow] (ode) -- (data);
\end{tikzpicture}
\caption{The continuous-time score view separates the known forward SDE from two valid reverse constructions: a stochastic reverse SDE and a deterministic probability-flow ODE.}
\label{fig:sde-ode}
\end{figure}

\begin{calloutmemoryhook}
Reverse diffusion needs the score; the same marginals can also be followed by a deterministic probability-flow ODE.
\end{calloutmemoryhook}

\section{Flow matching: learn the motion directly}

A continuous normalising flow generates data by integrating a time-dependent velocity field:
\[
\boxed{
\underbrace{\frac{dx_t}{dt}}_{\text{sample motion}}
=
\underbrace{v_{\theta}(x_t,t)}_{\text{learned velocity field}}
}.
\]
\textbf{Flow matching} trains this velocity field by first choosing a probability path between a simple source distribution and the data distribution, then regressing the velocity required along that path.
In its abstract form,
\[
\underbrace{\mathcal L_{\text{FM}}}_{\text{flow-matching loss}}
=
\mathbb E_{t,x_t}
\left[
\underbrace{\left\|v_\theta(x_t,t)-u_t(x_t)\right\|_2^2}_{\text{match the target velocity field}}
\right].
\]
The important design choice has therefore moved from a noise schedule and reverse stochastic process to the \textbf{probability path / coupling} and the velocity field that transports mass along it.

A useful special case pairs noise $x_{\text{noise}}$ with data $x_{\text{data}}$ and uses a straight interpolation:
\[
\underbrace{x_t}_{\text{point on the path}}
=
\underbrace{(1-t)x_{\text{noise}}+t x_{\text{data}}}_{\text{straight interpolation}},
\qquad
\underbrace{\frac{dx_t}{dt}}_{\text{target pairwise velocity}}
=
\underbrace{x_{\text{data}}-x_{\text{noise}}}_{\text{straight-line displacement}}.
\]
This is a useful intuition for rectified / straight-line flows, but general flow matching can use other probability paths and couplings.

The distinction from score modelling is conceptual:
\[
\underbrace{\nabla_x\log p_t(x)}_{\text{score}}
=\text{where density locally increases},
\qquad
\underbrace{v_t(x)}_{\text{velocity}}
=\text{where a sample should move next}.
\]
These fields can be mathematically related for particular paths, but they answer different local questions.

\begin{figure}[H]
\centering
\begin{tikzpicture}[>=Latex]
% left panel score
\begin{scope}[xshift=0cm]
\draw[rounded corners=2pt] (-3,-2.1) rectangle (3,2.2);
\node[font=\bfseries] at (0,1.85) {Score field};
\fill (-1.2,0) circle (2pt);
\fill (1.2,0) circle (2pt);
\draw[thick] (-1.2,0) ellipse (0.7 and 0.95);
\draw[thick] (1.2,0) ellipse (0.7 and 0.95);
\foreach \x/\y/\tx/\ty in {-2.4/1.2/-1.55/0.45,-2.4/-1.1/-1.55/-0.4,-0.2/1.25/-0.85/0.5,-0.2/-1.2/-0.85/-0.5,2.4/1.2/1.55/0.45,2.4/-1.1/1.55/-0.4,0.2/1.25/0.85/0.5,0.2/-1.2/0.85/-0.5}
{\draw[->, thick] (\x,\y) -- (\tx,\ty);}
\node[align=center] at (0,-1.75) {local density geometry};
\end{scope}
% right panel flow
\begin{scope}[xshift=7cm]
\draw[rounded corners=2pt] (-3,-2.1) rectangle (3,2.2);
\node[font=\bfseries] at (0,1.85) {Velocity field};
\fill (-2.2,-0.8) circle (1.7pt);
\fill (-2.2,0) circle (1.7pt);
\fill (-2.2,0.8) circle (1.7pt);
\fill (2.1,-0.7) circle (2pt);
\fill (2.1,0.7) circle (2pt);
\draw[->, thick] (-2.05,-0.8) .. controls (-0.8,-1.1) and (0.8,-0.9) .. (1.9,-0.7);
\draw[->, thick] (-2.05,0) .. controls (-0.5,0.1) and (0.5,0.6) .. (1.9,0.7);
\draw[->, thick] (-2.05,0.8) .. controls (-0.8,1.0) and (0.8,0.95) .. (1.9,0.7);
\node[align=center] at (0,-1.75) {transport through time};
\end{scope}
\end{tikzpicture}
\caption{Score and flow fields contain different local information. A score points toward higher current density; a velocity field specifies temporal motion along a chosen transport path.}
\label{fig:score-vs-flow}
\end{figure}

\begin{calloutspecialcase}
A straight interpolation with target velocity $x_{\text{data}}-x_{\text{noise}}$ is a useful rectified-flow mental model, not the definition of all flow matching.
\end{calloutspecialcase}

\begin{calloutmemoryhook}
Flow matching learns how samples should move.
\end{calloutmemoryhook}

\section{Classifier-free guidance: amplify the conditioning direction}

Conditional generation asks for samples from a distribution such as $p(x\mid c)$, where $c$ may be a class, text embedding or other context.
\textbf{Classifier-free guidance (CFG)} trains one model to operate both with the condition and with an intentionally dropped condition, so inference can compare the two predictions without a separate classifier.
For a generic model output $\hat y$,
\[
\boxed{
\underbrace{\hat y_{\text{guided}}}_{\text{prediction used by the sampler}}
=
\underbrace{\hat y_{\text{uncond}}}_{\text{what fits without the condition}}
+
\underbrace{s_{\text{guidance}}
\left(\hat y_{\text{cond}}-\hat y_{\text{uncond}}\right)}_{\text{amplified direction attributable to the condition}}
}.
\]
The output $\hat y$ may be a noise, score or velocity parameterisation; CFG acts on the model prediction used by the sampler rather than defining the underlying generative family.
At $s_{\text{guidance}}=1$ this expression returns the ordinary conditional prediction, while larger values extrapolate farther in the conditional direction.
Stronger guidance usually increases condition adherence but can reduce diversity and eventually distort or oversaturate samples.

\begin{figure}[H]
\centering
\begin{tikzpicture}[scale=1.15, >=Latex]
\coordinate (o) at (0,0);
\coordinate (u) at (2.2,0.7);
\coordinate (c) at (3.0,1.8);
\coordinate (g) at (4.6,3.25);
\draw[->, thick] (o) -- (u) node[midway, below] {$\hat y_{\text{uncond}}$};
\draw[->, thick] (o) -- (c) node[midway, above left] {$\hat y_{\text{cond}}$};
\draw[->, thick, dashed] (u) -- (c) node[midway, right] {$\hat y_{\text{cond}}-\hat y_{\text{uncond}}$};
\draw[->, very thick] (o) -- (g) node[midway, above left] {$\hat y_{\text{guided}}$};
\draw[densely dotted] (u) -- (g);
\fill (o) circle (1.5pt);
\node[below left] at (o) {origin};
\end{tikzpicture}
\caption{CFG treats the conditional-minus-unconditional difference as a guidance direction and scales it during inference. The diagram is geometric intuition, not a claim that model-output space is globally linear.}
\label{fig:cfg-vector}
\end{figure}

\begin{calloutdegreeofexperiment}
Vary conditioning-drop probability during training and guidance strength during sampling; measure condition adherence, diversity, sample quality and failure rate rather than assuming stronger guidance is universally better.
\end{calloutdegreeofexperiment}

\begin{calloutmemoryhook}
Classifier-free guidance amplifies what the condition changed.
\end{calloutmemoryhook}

\section{Modern composition: objective $\neq$ representation $\neq$ backbone $\neq$ sampler}

Modern image and multimodal systems commonly combine several ideas that live at different design layers.
A \textbf{latent diffusion / latent flow} system first compresses data with an autoencoder-like representation, learns diffusion or flow dynamics in that cheaper latent space, then decodes the generated latent back to the original space.
This trades some representation fidelity for much lower generative-model compute because the expensive field network operates on fewer spatial dimensions.
A U-Net or Transformer/DiT can then parameterise the denoising, score or velocity field, while CFG and the numerical sampler remain separate inference choices.

\begin{figure}[H]
\centering
\begin{tikzpicture}[
    >=Latex,
    node distance=9mm and 11mm,
    box/.style={draw, rounded corners=2pt, minimum width=23mm, minimum height=8mm, align=center, inner sep=3pt},
    arrow/.style={->, thick}
]
\node[box] (data) {data $x$};
\node[box, right=of data] (enc) {latent encoder};
\node[box, right=of enc] (latent) {latent $z$};
\node[box, right=of latent, minimum width=31mm] (field) {field network\\U-Net / DiT / other};
\node[box, right=of field] (zgen) {generated latent};
\node[box, right=of zgen] (dec) {latent decoder};
\node[box, right=of dec] (out) {generated $x$};
\node[box, above=9mm of field] (cond) {condition $c$};
\node[box, below=9mm of field] (sampler) {sampler / solver\\DDIM, SDE, ODE, \ldots};
\draw[arrow] (data) -- (enc);
\draw[arrow] (enc) -- (latent);
\draw[arrow] (latent) -- (field);
\draw[arrow] (field) -- (zgen);
\draw[arrow] (zgen) -- (dec);
\draw[arrow] (dec) -- (out);
\draw[arrow] (cond) -- (field);
\draw[arrow] (sampler) -- (field);
\end{tikzpicture}
\caption{A modern generative system composes independent design layers: representation, field-prediction backbone, conditioning and numerical sampler. Treating them as one method creates category errors.}
\label{fig:modern-generative-stack}
\end{figure}

\begin{calloutintuition}
``DiT'', ``latent diffusion'', ``flow matching'', ``CFG'' and ``DDIM'' are not five competing model families: they can describe different layers of one system.
\end{calloutintuition}

\section{Research evolution: which object should generation learn?}

The useful history is an accumulation of different answers to the same transport problem rather than a clean sequence in which one family replaces the previous one.

\begin{table}[H]
\centering
\small
\begin{tabularx}{\linewidth}{>{\bfseries}l X X X}
\toprule
Step & What changed & Problem targeted & New trade-off \\
\midrule
VAE & probabilistic latent space plus decoder & sampleable structured representation & reconstruction versus prior compatibility \\
GAN & adversarially learned generator & direct high-quality samples & unstable coupled optimisation / mode coverage \\
Score matching & learn local log-density gradients & avoid explicit normalised density modelling & sampling dynamics still required \\
DDPM & train denoising across known noise levels & stable scalable score-like learning & many reverse evaluations \\
DDIM / improved solvers & separate sampler from trained denoiser & reduce inference cost / change trajectory & numerical approximation and sampler choices \\
Score SDE & continuous-time diffusion formulation & unify score models and diffusion paths & additional continuous-time machinery \\
Flow matching & directly regress transport velocity & train ODE transport without simulating a reverse SDE & probability path / coupling becomes a major design choice \\
Distillation / consistency-style methods & compress many-step trajectories & sampling latency & extra training and approximation error \\
\bottomrule
\end{tabularx}
\caption{A compact evolution of generative mechanisms. Each step changes which object is learned or how the learned field is sampled.}
\label{tab:generative-evolution}
\end{table}

\section{Research optionality: what exactly can be changed?}

As in attention research, the useful question is not only ``which generative model is best?'' but \textbf{which design layer is being changed}.

\begin{table}[H]
\centering
\small
\begin{tabularx}{\linewidth}{>{\bfseries}l X X}
\toprule
Degree of experiment & Representative choices & Mechanism to measure \\
\midrule
Representation & pixels, compressed latent, discrete tokens & reconstruction loss, compute, memory, retained detail \\
Generative mechanism & adversarial, latent-variable, diffusion/score, flow & stability, mode coverage, likelihood/sample quality, sampling cost \\
Probability path & noise schedule, interpolation, transport coupling & target difficulty, path curvature, number of useful steps \\
Prediction target & $\epsilon$, $x_0$, score, $v$, velocity field & target scale across SNR/time, optimisation stability \\
Backbone & U-Net, CNN, Transformer/DiT, hybrid & quality versus parameters, FLOPs, memory \\
Conditioning route & concatenation, cross-attention, modulation & controllability, compute, information bottleneck \\
Guidance & none, classifier, classifier-free; vary strength & adherence versus diversity / artefacts \\
Time representation & discrete embedding, continuous embedding & interpolation across timesteps / solver compatibility \\
Sampler / solver & DDPM, DDIM, SDE solver, ODE solver & function evaluations, latency, stochasticity, error \\
Sampling budget & number and placement of steps & quality--latency curve \\
Training weighting & loss weight across time / SNR & which noise regimes dominate learning \\
\bottomrule
\end{tabularx}
\caption{Generative modelling has several distinct research layers; ablations should avoid conflating architectural, objective and solver changes.}
\label{tab:generative-optionality}
\end{table}

\begin{calloutdegreeofexperiment}
For diffusion/flow research, vary representation, probability path, prediction target, backbone, conditioning route, guidance, time weighting, solver and sampling budget; ablate one design layer at a time where possible.
\end{calloutdegreeofexperiment}

\section{Inductive bias and model choice}

Each family assumes that a particular intermediate learning problem is easier than modelling the data distribution directly.

\begin{table}[H]
\centering
\small
\begin{tabularx}{\linewidth}{>{\bfseries}l X X X}
\toprule
Method & Useful assumption & What the assumption buys & Main restriction / cost \\
\midrule
VAE & complex data can be organised through a simple latent prior & sampleable latent representation and interpolation & prior regularisation can discard reconstruction detail \\
GAN & a learned critic provides useful gradients for sample realism & direct one-pass generator after training & adversarial instability and possible mode collapse \\
Diffusion & a hard distribution can be reached by reversing many easy local corruptions & stable supervised denoising targets across noise levels & iterative inference \\
Score model & local log-density geometry is sufficient to build sampling dynamics & no need to model a normalised density directly & numerical dynamics still required \\
Flow matching & a useful probability path and coupling can be specified & direct velocity regression for ODE transport & a poor path can make the learned transport unnecessarily difficult \\
CFG & conditional--unconditional difference isolates useful conditioning information & inference-time controllability without a separate classifier & excessive guidance can reduce diversity and distort samples \\
\bottomrule
\end{tabularx}
\caption{Generative methods as inductive biases: each turns generation into a different intermediate prediction problem.}
\label{tab:generative-inductive-bias}
\end{table}

\section{Final compression}

If you remember only the core distinctions:
\begin{enumerate}
    \item A generative model learns a route from an easy source distribution to the data distribution.
    \item A VAE makes an autoencoder's latent space probabilistic and sampleable by balancing reconstruction against prior compatibility.
    \item A GAN learns a direct generator from a discriminator whose own parameters change during training.
    \item The score $\nabla_x\log p_t(x)$ is local log-density geometry; under Gaussian corruption it is directly related to the noise that diffusion models predict.
    \item Forward diffusion deliberately destroys information with a known process; the neural network learns reverse information.
    \item DDPM versus DDIM is mainly a sampler distinction: the same denoiser can support stochastic or deterministic-style trajectories.
    \item In continuous time, reverse diffusion uses the score inside a reverse SDE; a probability-flow ODE follows the same marginal distributions deterministically under an exact score.
    \item Flow matching directly learns a velocity field for a chosen probability path rather than first defining generation as reverse stochastic corruption.
    \item Classifier-free guidance scales the conditional-minus-unconditional model prediction during sampling.
    \item Representation, backbone, generative objective, conditioning, guidance and sampler are separate research layers and can be combined.
\end{enumerate}

\begin{calloutmemoryhook}
When reading a new generative-model paper, ask: what distribution path is assumed, what field/object is learned, what target teaches it, and how is that object integrated at inference?
\end{calloutmemoryhook}